\chapter{Tinjauan Pustaka}

\section{Penelitian Terdahulu}
\TemplateTodo{Ringkas penelitian terdahulu yang paling relevan. Jelaskan tujuan, metode, data, hasil, dan keterbatasannya. Akhiri dengan posisi penelitian Anda dibandingkan karya-karya tersebut.}

Tinjauan pustaka bukan sekadar daftar ringkasan artikel. Susun pembahasan secara tematik sehingga pembaca memahami perkembangan pengetahuan, kesenjangan yang masih ada, dan alasan penelitian ini diperlukan.

\begin{table}[H]
\centering
\caption{Contoh format ringkasan penelitian terdahulu}
\begin{tabularx}{\linewidth}{@{}p{0.18\linewidth}X X@{}}
\toprule
Peneliti & Fokus dan Metode & Keterkaitan dengan \DocTypeName \\
\midrule
Nama peneliti (tahun) & Ringkasan pendek metode, data, dan hasil utama. & Jelaskan pelajaran, celah, atau pembeda penelitian Anda. \\
\bottomrule
\end{tabularx}
\end{table}

\section{Landasan Teori}
\TemplateTodo{Jelaskan konsep, teori, algoritma, standar, atau kerangka kerja yang benar-benar digunakan pada penelitian.}

\subsection{Konsep Utama Pertama}
Tuliskan definisi, prinsip kerja, dan relevansi konsep terhadap penelitian. Gunakan sitasi untuk definisi atau teori yang bukan pendapat pribadi.

\subsection{Konsep Utama Kedua}
Tambahkan konsep lain sesuai kebutuhan. Jika menggunakan rumus, jelaskan arti setiap simbol dan alasan rumus tersebut digunakan.

\section{Kerangka Pemikiran}
\TemplateTodo{Jelaskan hubungan antara masalah, teori, metode, data, dan keluaran penelitian. Bagian ini dapat dibantu dengan gambar alur konseptual.}

\begin{figure}[H]
\centering
\includegraphics[width=0.75\linewidth]{example-figure.pdf}
\caption{Contoh penggunaan gambar pada template}
\label{fig:example-figure}
\end{figure}

Contoh penggunaan gambar ditunjukkan pada Gambar~\ref{fig:example-figure}. Notasi matematika dapat digunakan untuk menjelaskan hubungan antar lapisan pada Gambar~\ref{fig:example-figure}. Misalkan vektor masukan dinyatakan sebagai
\begin{equation}
\mathbf{x} =
\begin{bmatrix}
x_1 & x_2 & \cdots & x_n
\end{bmatrix}^{\mathsf{T}}.
\end{equation}
Untuk lapisan tersembunyi ke-$l$, aktivasi neuron dihitung dengan
\begin{equation}
\mathbf{a}^{[l]} = \sigma\left(\mathbf{W}^{[l]}\mathbf{a}^{[l-1]} + \mathbf{b}^{[l]}\right),
\qquad l = 1,2,\ldots,L-1,
\end{equation}
dengan $\mathbf{a}^{[0]}=\mathbf{x}$, $\mathbf{W}^{[l]}$
adalah matriks bobot pada lapisan ke-$l$,
$\mathbf{b}^{[l]}$ adalah vektor bias, dan
$\sigma(\cdot)$ adalah fungsi aktivasi. Keluaran model dinyatakan sebagai
\begin{equation}
\hat{\mathbf{y}} = f\left(\mathbf{W}^{[L]}\mathbf{a}^{[L-1]} + \mathbf{b}^{[L]}\right),
\end{equation}
dengan komponen keluaran
\begin{equation}
\hat{\mathbf{y}} =
\begin{bmatrix}
\hat{y}_1 & \hat{y}_2 & \cdots & \hat{y}_k
\end{bmatrix}^{\mathsf{T}}.
\end{equation}
